% A deliberately plain paper that a stock LaTeX install can build.
%
% The other LaTeX fixture uses a conference class, natbib, and a figure file
% that does not exist, because it exists to exercise ingest against what real
% papers look like. This one exists to prove that a paper still compiles after
% the ledger has been applied to it, which means it has to compile in the first
% place, on texlive-latex-base and nothing else.
%
% Three things are planted on purpose. Filler, so the ledger is not empty.
% Sentences of very different lengths, so the author's length band is wide
% enough that the voice lock does not refuse every cut and leave the compile
% test proving nothing. And filler sitting right beside inline math and a
% cross-reference, so an applied cut has to miss a protected range by inches
% rather than by half a page.
\documentclass{article}

\title{Sparse Retrieval at Equal Cost}
\author{A. Researcher}

\begin{document}
\maketitle

\section{Introduction}
\label{sec:intro}

It is important to note that dense encoders are reported to beat sparse
baselines on this benchmark. We present a comparison that holds the budget
fixed rather than holding the index size fixed, and we report recall at ten
for every configuration we measured in this study, across the three corpora we
had access to and across the five random seeds we ran each configuration on.
We tested three systems.

Needless to say, the gap closes once the cost is held fixed.

\section{Method}

It should be mentioned that the corpus is indexed with BM25 using
$k_1 = 0.9$ and $b = 0.4$.
Queries are expanded with the top three terms taken from a first pass over the
corpus, and the expansion is applied before the second pass rather than after
it, which matters because the two orders give different results on the shorter
queries in this benchmark.
We ran each pass once.

\section{Results}

Recall at ten rises from 0.62 to 0.71 with expansion enabled, and the cost is
one third of the dense baseline over the five thousand queries we sampled for
this comparison, measured on the same hardware for every system we compared.
We significantly outperform the dense encoder at the budget we held fixed in
Section~\ref{sec:intro}.
The gain holds on every corpus.

\end{document}
